%% Genomics Intelligence Architecture — Conditional Program Model
%% Migrated from singlet-intelligence/docs/architecture/cpm.md
\documentclass[10pt,twocolumn,letterpaper]{article}
\usepackage{../../templates/singletai-preprint}
\usepackage{graphicx}
\usepackage{amsmath,amssymb}
\usepackage{booktabs}
\usepackage{algorithm}
\usepackage{algpseudocode}
\usepackage{natbib}
\bibliographystyle{unsrtnat}

% ── Metadata ─────────────────────────────────────────────────────
\title{The Conditional Program Model: A Generative Architecture \\
  for Transcriptomics Intelligence}
\author{
  Zach DeBruine\textsuperscript{1,*}
  \\[6pt]
  \textsuperscript{1}Department of Computing, Grand Valley State University, Allendale, MI 49401
}

\begin{document}
\maketitle
\correspondingauthor{debruinz@gvsu.edu}

\begin{abstract}
Single-cell transcriptomics generates expression profiles for millions of cells,
yet researchers can only query conditions that have been experimentally observed.
We propose the Conditional Program Model (CPM), a generative architecture that
extends non-negative matrix factorization (NMF) with a conditional neural
encoder to predict gene program activities for arbitrary, unobserved biological
conditions. The CPM consists of three components: (1)~a gene program dictionary
$W \in \mathbb{R}_{\geq 0}^{g \times k}$ learned via regularized NMF on an
atlas of ${\sim}10^9$ cells, with angular regularization enabling $k = 1024$--$2048$
programs without degeneracy; (2)~a conditional encoder $f_\theta$ mapping
categorical and continuous metadata to program activity distributions via
cross-attention over learned embeddings; and (3)~a parameter-free decoder
$\hat{x} = W\hat{h}$ that translates predictions to gene space. The encoder
predicts both mean $\mu$ and variance $\sigma$ of program activities, capturing
biological variability across donors and cell states. We describe extensions for
cross-species integration via ortholog-projected shared program spaces,
perturbation prediction through learned intervention operators in program space,
and natural language hypothesis testing by parsing biological queries into
structured model inference calls. The architecture separates the 45~GB gene
program dictionary (server-side) from the 12~MB encoder (client-side), enabling
offline inference on commodity hardware. We present a validation framework
including reconstruction accuracy ($r > 0.85$ per cell type), cross-validation
of rank selection via speckled holdout, and zero-shot condition prediction
benchmarks.
\end{abstract}

\keywords{single-cell transcriptomics \and non-negative matrix factorization
  \and generative model \and gene programs \and conditional inference}

% ── Introduction ─────────────────────────────────────────────────
\section{Introduction}
\label{sec:intro}

The single-cell transcriptomics community currently answers biological questions
by operating directly on observed data---downloading count matrices, subsetting by
condition, computing differential expression, running enrichment analyses, and
fitting dimensionality reductions. This workflow has two structural limitations.

First, one can only answer questions about conditions that have been
experimentally observed. If a researcher wants to know what happens when gene~$X$
is knocked out in cell type~$Y$ under disease condition~$Z$, they need a dataset
where someone actually performed that experiment. Most condition combinations
have never been observed.

Second, the bottleneck is data wrangling, not analysis. Researchers spend
60--80\% of their time finding, downloading, formatting, and quality-controlling
data~\citep{zheng2017massively}. The actual biological insight---``which programs
changed?''---takes minutes once the data is prepared.

A transcriptomics intelligence platform inverts this paradigm. Instead of
operating on data, it uses all observed data as training substrate to learn a
model that can: (i)~\emph{generate} predicted expression profiles for arbitrary
condition combinations; (ii)~\emph{decompose} those predictions into named,
interpretable biological programs; (iii)~\emph{predict} how perturbations shift
program activities; (iv)~\emph{transfer} knowledge across species via shared
program structure; and (v)~\emph{answer questions in natural language} by parsing
queries into structured model inference calls.

The key question is: what model architecture achieves all of this while remaining
interpretable, scalable, and grounded in published biology?

\subsection{NMF as a foundation, not a generative model}

Non-negative matrix factorization decomposes a gene expression matrix
$X \in \mathbb{R}_{\geq 0}^{g \times n}$ into
\begin{equation}
  X \approx WH
  \label{eq:nmf}
\end{equation}
where $W \in \mathbb{R}_{\geq 0}^{g \times k}$ contains $k$ gene program
loadings and $H \in \mathbb{R}_{\geq 0}^{k \times n}$ contains program activity
levels for each cell~\citep{lee1999learning,debruine2021rcppml}.

Given only~$W$ (the learned gene programs), NMF has no mechanism to produce~$h$
(the program activity vector) for a new, unobserved condition. NMF can decompose
an observed expression vector---given~$x$, solve
$\min_{h \geq 0} \|x - Wh\|^2$---and reconstruct from known activities---given~$h$,
compute $\hat{x} = Wh$. But it \emph{cannot} answer: ``What would~$h$ be for a
human hepatocyte under NASH with a PCSK9 knockout?'' without first observing that
cell.

What NMF does provide is the critical ingredient that other architectures lack:
an interpretable, biologically grounded output space. Each dimension of~$h$
corresponds to a named biological program (e.g., oxidative phosphorylation, T~cell
activation, cell cycle G2/M). This interpretability is not a nicety---it is the
product. A black-box expression prediction is scientifically useless; a
prediction decomposed into ``program~12 (lipid metabolism) increases 2.3$\times$''
is actionable.

The architecture must therefore \emph{not replace} NMF but \emph{extend it} with
a conditional generative component that maps metadata to program activities,
while preserving NMF's gene programs as the structured, interpretable output
vocabulary.

% ── Methods ──────────────────────────────────────────────────────
\section{Architecture: The Conditional Program Model}
\label{sec:cpm}

The Conditional Program Model has three components (Figure~\ref{fig:architecture}):
\begin{enumerate}
  \item \textbf{Gene Program Dictionary}~$W$---learned once via NMF on the full
    atlas, fixed during inference.
  \item \textbf{Conditional Program Encoder}~$f_\theta$---a neural network
    mapping metadata conditions to program activity distributions.
  \item \textbf{Expression Decoder}---the matrix multiplication
    $\hat{x} = W\hat{h}$, which is parameter-free.
\end{enumerate}

The generation pathway is:
\begin{equation}
  \text{metadata} \xrightarrow{f_\theta} \hat{h} \xrightarrow{\times W} \hat{x}
  \label{eq:pathway}
\end{equation}

\subsection{Gene program dictionary via regularized NMF}
\label{sec:nmf}

The gene program dictionary is learned via NMF on the full atlas
($g \approx 30{,}000$ genes, $n \approx 10^9$ cells). The objective is:
\begin{equation}
  \min_{W, H \geq 0} \|X - WH\|_F^2
    + \lambda_{\ell_1}(\|W\|_1 + \|H\|_1)
    + \lambda_{\text{ang}} \!\!\sum_{j \neq j'}\!\!
      \max\!\bigl(0, \cos(W_{:,j}, W_{:,j'}) - \epsilon\bigr)^2
  \label{eq:nmf-obj}
\end{equation}

The $\ell_1$ sparsity penalty promotes sparse gene loadings (each program driven
by a focused gene set) and sparse activities (each cell uses a subset of
programs), tuned so that the Gini coefficient of each cell's $h_i$ vector exceeds
0.7. The angular regularization penalty penalizes program columns whose cosine
similarity exceeds threshold~$\epsilon$, promoting orthogonality between
programs. With angular regularization, ranks of 1024--2048 become recoverable
without degeneracy; without it, programs at high ranks collapse into
near-duplicates.

The number of programs~$k$ is determined via speckled holdout cross-validation:
randomly mask 1\% of nonzero entries in~$X$, train NMF on remaining entries,
evaluate reconstruction error on held-out entries, and select the~$k$ minimizing
holdout error before overfitting.

\subsection{Conditional program encoder}
\label{sec:encoder}

The encoder is a neural network $f_\theta : \mathcal{M} \to
\mathbb{R}_{\geq 0}^k$ mapping metadata space~$\mathcal{M}$ to program activity
distributions. Each cell's metadata includes:

\begin{table}[t]
  \centering
  \caption{Metadata fields for the conditional encoder.}
  \label{tab:metadata}
  \small
  \begin{tabular}{@{}lll@{}}
    \toprule
    Field & Type & Representation \\
    \midrule
    Tissue & Categorical & Learned embedding \\
    Cell type & Categorical & Learned embedding \\
    Disease & Categorical & Learned embedding \\
    Developmental stage & Continuous & Scalar (days) + embedding \\
    Sex & Categorical & Learned embedding \\
    Perturbation & Set & Gene embedding(s) \\
    Assay & Categorical & Learned embedding \\
    Donor & Categorical & Learned embedding \\
    Cell cycle phase & Categorical & Learned embedding \\
    \bottomrule
  \end{tabular}
\end{table}

Categorical fields receive learned embeddings initialized randomly and trained
end-to-end. Continuous scalars (developmental stage, doublet score, mitochondrial
fraction) pass through a small MLP to produce $d$-dimensional vectors that
concatenate with categorical embeddings. Missing fields are replaced with a
learned \texttt{[UNKNOWN]} token, and the model is trained with random metadata
dropout for robustness to partial information.

\paragraph{Cross-attention.}
All embeddings attend to each other via a cross-attention layer that captures
interactions (e.g., tissue$\times$disease, cell\_type$\times$developmental\_time).
The attended representations are concatenated and passed through an MLP trunk:
\begin{align}
  \mathbf{z} &= \text{CrossAttn}(\mathbf{e}_1, \ldots, \mathbf{e}_m) \\
  \mu &= \text{Softplus}\bigl(\text{MLP}_\mu(\mathbf{z})\bigr) \\
  \sigma &= \text{Softplus}\bigl(\text{MLP}_\sigma(\mathbf{z})\bigr)
\end{align}

The encoder predicts both mean~$\mu$ and variance~$\sigma$ of program
activities, capturing biological variability. The softplus activation ensures
non-negativity.

\subsection{Training objective}
\label{sec:training}

For each cell~$i$ with NMF-derived activities~$h_i$ and metadata~$m_i$, the
primary loss is Gaussian negative log-likelihood:
\begin{equation}
  \mathcal{L}_{\text{NLL}} = \sum_{i=1}^{n} \sum_{j=1}^{k}
    \left[
      \frac{(h_{ij} - \mu_j(m_i))^2}{2\sigma_j(m_i)^2} + \log \sigma_j(m_i)
    \right]
  \label{eq:nll}
\end{equation}

The total objective includes program correlation, reconstruction, and
calibration regularizers:
\begin{equation}
  \mathcal{L} = \mathcal{L}_{\text{NLL}}
    + \lambda_{\text{corr}} \mathcal{L}_{\text{corr}}
    + \lambda_{\text{recon}} \mathcal{L}_{\text{recon}}
    + \lambda_{\text{cal}} \mathcal{L}_{\text{cal}}
  \label{eq:total-loss}
\end{equation}

% ── Cross-Species ────────────────────────────────────────────────
\section{Cross-Species Integration}
\label{sec:cross-species}

Cross-species transfer exploits the structural conservation of gene programs
across species. Given species $A$ and $B$ with ortholog mapping matrix
$O \in \{0, 1\}^{g_A \times g_B}$, we project species~$B$'s program dictionary
into species~$A$'s gene space:
\begin{equation}
  W_B^{(A)} = O^\top W_B
  \label{eq:ortholog-projection}
\end{equation}

Programs are aligned across species via bipartite matching on cosine similarity
between $W_A$ and $W_B^{(A)}$ columns. Programs with cosine similarity exceeding
0.4 are designated as \emph{shared programs}; species-specific residual programs
capture lineage-specific biology.

This enables zero-shot cross-species prediction: given human metadata, generate
mouse-space expression by using shared program activities with the mouse gene
program dictionary.

% ── Perturbation ─────────────────────────────────────────────────
\section{Perturbation Prediction}
\label{sec:perturbation}

Perturbation prediction operates in program space rather than gene space.
A perturbation encoder maps the intervention (gene knockout, drug treatment,
overexpression) to a perturbation operator:
\begin{equation}
  \hat{h}_{\text{perturbed}} = \hat{h}_{\text{baseline}} + \Delta h(\text{perturbation})
  \label{eq:perturbation}
\end{equation}

where $\Delta h$ is predicted by a perturbation-specific network that takes
the baseline program activities and perturbation description as input. Operating
in the $k$-dimensional program space (rather than the $g$-dimensional gene space)
makes perturbation prediction tractable and interpretable: the model predicts
which biological programs change and by how much.

% ── Deployment ───────────────────────────────────────────────────
\section{Deployment Architecture}
\label{sec:deployment}

The CPM separates into server and client components:

\begin{table}[t]
  \centering
  \caption{Server--client architecture for CPM deployment.}
  \label{tab:deployment}
  \small
  \begin{tabular}{@{}lll@{}}
    \toprule
    Component & Size & Location \\
    \midrule
    Gene program dictionary $W$ & $\sim$45~GB & Server (GPU) \\
    Conditional encoder $f_\theta$ & $\sim$12~MB & Client (CPU) \\
    Program metadata/annotations & $\sim$2~MB & Client \\
    \bottomrule
  \end{tabular}
\end{table}

Client-side inference requires only the encoder and program annotations.
Full gene-level reconstruction requires server access for the $W$ matrix
multiplication. The \texttt{singlet} Python client handles this split
transparently via the SingletDB API.

% ── Validation ───────────────────────────────────────────────────
\section{Validation Framework}
\label{sec:validation}

We propose a multi-level validation framework:

\paragraph{Level 1: Reconstruction accuracy.}
For each cell type with $\geq$1000 cells, split 80/20 into train/test. Train
NMF + encoder on training cells, predict on test cell metadata, compare
$\hat{x}$ vs.\ observed~$x$. Target: Pearson $r > 0.85$ per-cell-type mean
profile.

\paragraph{Level 2: Zero-shot condition prediction.}
Hold out entire conditions (e.g., all cells under disease~$D$ in tissue~$T$)
from training. Predict program activities for the held-out condition using the
encoder trained on remaining data. This tests genuine generalization.

\paragraph{Level 3: Perturbation benchmarks.}
Evaluate on published Perturb-seq datasets: given control cell profiles and
perturbation identity, predict post-perturbation profiles. Compare against
GEARS, scGen, and CPA on standard metrics (Pearson $r$, top-20 differential
gene overlap).

\paragraph{Level 4: Cross-species transfer.}
Train on human atlas, predict mouse cell type profiles via ortholog projection.
Evaluate on held-out mouse cell types against observed mouse data.

% ── Discussion ───────────────────────────────────────────────────
\section{Discussion}
\label{sec:discussion}

The CPM architecture addresses a fundamental gap in single-cell analysis: the
inability to query conditions that have not been experimentally observed. By
extending NMF with a conditional encoder, the architecture preserves the
interpretability of gene programs while enabling generative inference.

Several design choices merit discussion. The separation of NMF training (Stage~1)
from encoder training (Stage~2) means the gene program dictionary is fixed during
conditional learning. This is a deliberate tradeoff: it ensures that program
meanings remain stable across all predictions, at the cost of not jointly
optimizing programs and predictions.

The choice of learned embeddings over ontology-initialized embeddings is
motivated by the observation that ontological distance does not always predict
transcriptomic similarity. With ${\sim}500$ cell types and ${\sim}200$ tissues,
the embedding spaces are small enough for the model to discover relationships
from data.

\paragraph{Limitations.}
The architecture assumes that gene programs are condition-invariant---that the
same biological processes operate across all conditions, only varying in activity.
This may not hold for highly aberrant states (e.g., cancer) where novel programs
emerge. The hierarchical NMF extension (Section~\ref{sec:nmf}) partially
addresses this.

% ── Conclusions ──────────────────────────────────────────────────
\section{Conclusions}
\label{sec:conclusions}

We have described the Conditional Program Model, a generative architecture for
transcriptomics intelligence that extends NMF with conditional inference. The
CPM decomposes prediction into an interpretable program vocabulary, supports
cross-species transfer and perturbation prediction, and separates into a compact
client-side encoder for offline use. This architecture enables a shift from
data-centric to model-centric analysis in single-cell genomics.

% ── Acknowledgments ──────────────────────────────────────────────
\section*{Acknowledgments}
This work was supported by Grand Valley State University.

\section*{Data \& Code Availability}
Source code: \url{https://github.com/Singlet-AI/singlet}.
Gene programs: \url{https://singletdb.com}.

% ── References ───────────────────────────────────────────────────
\bibliography{refs}

\end{document}
